%%%%%%%%%%%%%%%%%%%%%%%%%%%%%%%%%%%%%%%%%%%%%%%%%%%%%%%%%%%%%%%%%%%%%%%%
%                                                                      %
% LaTeX, FIIW thesis template                                          %
% 28/11/2014 v1.2                                                      %
%                                                                      %
%%%%%%%%%%%%%%%%%%%%%%%%%%%%%%%%%%%%%%%%%%%%%%%%%%%%%%%%%%%%%%%%%%%%%%%%
%\documentclass[11pt,a4paper]{report}
% Indien je je thesis recto-verso wil afdrukken gebruik je onderstaande opties i.p.v. bovenstaande
\documentclass[11pt,a4paper,twoside,openright]{report}

\usepackage[a4paper,left=3.5cm, right=2.5cm, top=3.5cm, bottom=3.5cm]{geometry}
\usepackage[dutch]{babel}
\usepackage{graphicx}
\usepackage[latin1]{inputenc}           % om niet ascii karakters rechtstreeks te kunnen inputten
%\usepackage[utf8]{inputenc}            % commentarieer deze regel uit als je utf8 encoded files gebruikt in plaats van latin1
\usepackage[square,numbers]{natbib}		% use [1] notation
\usepackage{listings}             		% voor het weergeven van broncode
\usepackage{verbatim}					% weergeven van code, commando's, ...
\usepackage{hyperref}					% maak PDF van de thesis navigeerbaar
\usepackage{url}						% URL's invoegen in tekst met behulp van \url{http://}
\usepackage[small,bf,hang]{caption}     % om de captions wat te verbeteren
\usepackage[final]{pdfpages}            % gebruikt voor het invoegen van het artikel in pdf-formaat
\usepackage{pslatex}					% andere lettertype's dan de standaard types
\bibliographystyle{IEEEtranN} 			% how formatted in bibliography
\usepackage{sectsty}					% aanpassen van de fonts van sections en chapters
\allsectionsfont{\sffamily}
\chapterfont{\raggedleft\sffamily}
\usepackage{xcolor}						% om todo's aan te maken
\newcommand\todo[1]{\textcolor{red}{#1}}
					% andere lettertype's dan de standaard types

\usepackage[acronym]{glossaries}        % voor afkortingen
\makeglossaries
\newacronym[plural=LPWANs,firstplural=Low-Power Wide-Area Networks (LPWANs)]{lpwan}{LPWAN}{Low-Power
	Wide-Area Network}
\newacronym{iot}{IoT}{Internet-of-Things} % default expression for defining a new acronym
\newacronym{ook}{OOK}{On-Off Keying}
\newacronym{ppm}{PPM}{Pulse Position Modulation}
\newacronym{ofdm}{OFDM}{Orthogonal Frequency Division Multiplexing}
\newacronym{csk}{CSK}{Color Shift Keying}
\newacronym{cap}{CAP}{Carrier-less Ampltitude and Phase}
\newacronym{aoa}{AOA}{Angle Of Arrival}
\newacronym{toa}{TOA}{Time Of Arrival}
\newacronym{rss}{RSS}{Received Signal Strength}
\newacronym[plural=LEDs,firstplural=Light Emitting Diodes (LEDs)]{led}{LED}{Light Emitting Diode}
\newacronym{vlc}{VLC}{Visible Light Communication}
\newacronym{vlp}{VLP}{Visible Light based Positioning}
\newacronym{rfid}{RFID}{Radio-Frequency Identification}
\newacronym{uwb}{UWB}{Ultra-WideBand}
\newacronym{ap}{AP}{Acces Point}
\newacronym{ask}{ASK}{Amplitude-Shift Keying}
\newacronym{nrz}{NRZ}{Non-return-to-zero}
\newacronym{rz}{RZ}{Return-to-zero}
\newacronym{vppm}{VPPM}{Variable PPM}
\newacronym{oppm}{OPPM}{Overlapping PPM}
\newacronym{mppm}{MPPM}{Multi-Pulse PPM}
\newacronym{isi}{ISI}{Inter-Symbol Interference}
\newacronym{ips}{IPS}{Indoor Positioning System}
\newacronym{pd}{PD}{Photodiode}
\newacronym{tdoa}{TDOA}{Time Difference of Arrival}
\newacronym{sbc}{SBC}{Single Board Computers}
\newacronym{ml}{ML}{Machine Learning}
\newacronym{ann}{ANN}{Artifici\"ele Neurale Netwerken}
\newacronym{nn}{NN}{Neurale Netwerken}
\newacronym{gpu}{GPU}{Graphics Processing Unit}
\newacronym{bbai}{BB AI}{BeagleBone Ai}
\newacronym{soc}{SoC}{System On Chip}
\newacronym{ai}{AI}{Artifici\"ele Intelligentie}
\newacronym{eve}{EVE}{embedded-vision-engine}
\newacronym{svm}{SVMs}{Support Vector Machines}
\newacronym{tpu}{TPU}{Tensor Processing Unit}
\newacronym{asic}{ASIC}{Application Specific Integrated Circuit}
\newacronym{som}{SoM}{System-on-Module}
\newacronym{tops}{TOPS}{tera operations per second}
\newacronym{cpu}{CPU}{Central Processing Unit}
\newacronym{relu}{ReLU}{rectified linear unit}



\usepackage{sectsty}					% aanpassen van de fonts van sections en chapters
\allsectionsfont{\sffamily}
\chapterfont{\raggedleft\sffamily}

\usepackage{float}                      % De optie H voor de plaatsing van figuren op de plaats waar je ze invoegt. bvb. \begin{figure}[H]
\usepackage{longtable}					% tabellen die over meerdere pagina's gespreid worden
%\usepackage[times]{quotchap}           % indien je fancy hoofdstuktitels wil
\usepackage[none]{hyphenat}
\usepackage{latexsym}
\usepackage{amsmath}
\usepackage{amssymb}

\usepackage{fiiw_gent}
% \usepackage{fiiw_ghent_eng} % For the english version (also change last page at the bottom of this file!

%door onderstaande regels in commentaar te zetten, of op false, kan je pagina's weglaten
%bijvoorbeeld het weglaten van een voorwoord, lijst met symbolen, ...
%%%%%%%%%%%%%%%%%%%%%%%%%%%%%%%%%%%%%%%%%%%%%%%%%%%%%%%%%%%%%%%%%%%%%%%%%%%%%%%%%%%%%%%%
%voorwoord toevoegen?
\acknowledgementspagetrue
\acknowledgements{voorwoord}			%.tex file met daarin het voorwoord
%abstract toevoegen?
\abstractpagetrue
\abstracts{abstract}					%.tex file met daarin het abstract
%lijst van figuren toevoegen?
\listoffigurespagetrue
%lijst van tabellen toevoegen?
\listoftablespagetrue
%lijst van symbolen toevoegen?
\listofsymbolspagetrue
\listofsymbols{symbolen}				%.tex file met daarin de lijst van symbolen



%informatie over het eindwerk, de promotor, ...
%%%%%%%%%%%%%%%%%%%%%%%%%%%%%%%%%%%%%%%%%%%%%%%
\opleiding{Industri\"ele Ingenieurswetenschappen Elektronica-ICT}
\afdeling{Embedded Systems}

%\campus{gent}

\title{Accelerated edge computing\\ with Neural Networks}

% \author{naam student}
\forename{Arno}
\surname{Plaetinck}
\academicyear{2019 - 2020}

\promotorA[Promotor]{Prof. Dr. Ir. Nobby Stevens}
\promotorB[Co-promotoren]{Ing. Willem Raes\\Ing. Jorik De Bruycker}

\begin{document}
\selectlanguage{dutch}

\preface

%second chapter of your thesis
\chapter{Literatuurstudie}
%intro
\section{Visible Light Communication}
	%RF bandwidths are failing to meet requirements, VLC is promising alternatives
	%korte intro
\section{Visible Light based Positioning}
	\subsection{techniques}
	%different technologies
		\subsubsection{LED Technologie}
			\begin{itemize}
				\item White LED's
				\item RGB-LED's
			\end{itemize}
		\subsubsection{Modulatie methode}
			\begin{itemize}
				\item On-Off Keying (OOK)
				\item Pulse Postion Modulation (PPM)
				\item Orthogonal Frequency Division Multiplexing (OFDM)
				\item Color Shift Keying (CSK)
				\item Carrier-less Amplitude and Phase (CAP)
			\end{itemize}
		\subsubsection{Type ontvanger}
			\begin{itemize}
				\item Photodiode
				\item Image Sensors
			\end{itemize}
	\subsection{Taxonomy}
		\subsubsection{Wiskundige methode}
			\begin{itemize}
				\item Proximity
				\item Triangulation
					\begin{itemize}
						\item AOA
						\item TOA
						\item RSS
					\end{itemize}
				\item Fingerprint
					\begin{itemize}
						\item Map-Based fingerprint
						\item Online stage/runtime stage
					\end{itemize}
			\end{itemize}
		\subsubsection{Sensor assisted methode}
		\subsubsection{Optimization methode}
\section{Artifici\"ele Neurale Netwerken}

% waarom op de edge plaatsen van EBC (Voor en Nadelen) (miss meer situring)
% Specs van de verschillende boards 
%Het eerste hoofdstuk van je thesis.
\chapter{Inleiding}
In dit hoofdstuk\index{hoofdstuk} gaan we een voorbeeld geven van een voetnoot\footnote{Dit is dus een voetnoot}. Een referentie naar hoofdstuk ~\ref{verwijzing}, dat zich op pagina \pageref{verwijzing} bevindt, is dus ook een koud kunstje. Zorg er wel voor dat je de namen van de labels een beetje verstandig kiest. Hoofdstukken label je het best als hfdstk:naam, plaatjes als img:naam en tabellen\index{tabellen} als tabel:naam. Zo verlies je zelf de bomen in het bos niet.
\newpage
SDffjfhd fsffh hsf
fh fhf
shf klfh
ffffsdfklfhklfhklfhhfklfhkldhffhsdfhfhfhfhfh
\newpage
dhfhffh hf fh fh fhfh fhfh hfh fhffhsdfhfhfhfhfhfhsdfh hfh fh
 




%second chapter of your thesis
\chapter{Bespreking}
In het vorig hoofdstuk hebben we naar deze tekst verwezen\label{verwijzing}.

%%%%%%%%%%%%%%%%%%%%%%%%%%%%%%%%%%%%%%%%%%%%%%%%%%%%%%%%%%%%%%%%%%% 
%                                                                 %
%                            CHAPTER                              %
%                                                                 %
%%%%%%%%%%%%%%%%%%%%%%%%%%%%%%%%%%%%%%%%%%%%%%%%%%%%%%%%%%%%%%%%%%% 
 
\chapter{This is the another Chapter}
 
You can say great work has been done about something \citep{Castleman98,Granlund95} or say that \citet{Holmes95} did something really great.
xxxx xxxxx xxxx xxxxxxxxx 
xxx xxxxx xxxxx xxx xxxx xxxx xxxxx xxxxx xxxx xxxxx xxxx xxxxxxxxx
 
\begin{figure}
\vspace{2.0in}
\caption{This is the Caption for Figure 1}
\end{figure}
 
xxx xxxxx xxxxx xxx xxxx xxxx xxxxx xxxxx xxxx xxxxx xxxx xxxxxxxxx
xxx xxxxx xxxxx xxx xxxx xxxx xxxxx xxxxx xxxx xxxxx xxxx xxxxxxxxx
xxx xxxxx xxxxx xxx xxxx xxxx xxxxx xxxxx xxxx xxxxx xxxx xxxxxxxxx
 
\begin{table}[top]
\begin{center}
\begin{tabular}{lll}
Here's       & an          & example  \\
of           & a           & table    \\
floated      & with        & the      \\
\verb+table+ & environment & command.
\end{tabular}
\end{center}
\caption{This is the Caption for Table 1}
\end{table}
 
xxx xxxxx xxxxx xxx xxxx xxxx xxxxx xxxxx xxxx xxxxx xxxx xxxxxxxxx
xxx xxxxx xxxxx xxx xxxx xxxx xxxxx xxxxx xxxx xxxxx xxxx xxxxxxxxx
 
\section{This is a Section Heading}
 
% Bibliografie: referenties. De items zitten in bibliografie.bib
%%%%%%%%%%%%%%%%%%%%%%%%%%%%%%%%%%%%%%%%%%%%%%%%%%%%%%%%%%%%%%%%%
% Indien je ook de niet geciteerde werken in je bibliografie wil opnemen, commentarieer dan onderstaande regel uit!
%\nocite{*}
\bibliographystyle{apalike}
\bibliography{bibliografie}

% Eventueel enkele appendices
%%%%%%%%%%%%%%%%%%%%%%%%%%%%%%
\appendix
\chapter{Een aanhangsel}
sdfsffqsfsf

% Bijlage met daarin het wetenschappelijk artikel
%%%%%%%%%%%%%%%%%%%%%%%%%%%%%%%%%%%%%%%%%%%%%%%%%%
\chapter{Beschrijving van deze masterproef in de vorm van een wetenschappelijk artikel}
%\includepdf{artikel.pdf}

% Bijlage met daarin de poster
%%%%%%%%%%%%%%%%%%%%%%%%%%%%%%%
\chapter{Poster}
%\includepdf{poster.pdf}

\includepdf{back_fiiw_gent.pdf}

\end{document}
